\chapter{Coma}

\section*{Definition}
Coma is a state of profound unconsciousness in which you cannot be aroused and lacks normal responses to external stimuli. It represents a medical emergency requiring immediate evaluation. Coma results from bilateral cerebral hemisphere dysfunction or brainstem reticular activating system disruption.

\section*{When to See a Doctor}
\begin{itemize}
\item Unresponsiveness to verbal or painful stimuli
 \item Any sudden or unexplained unresponsiveness
\item Bilateral pupillary abnormalities
\item Abnormal breathing patterns
\item Posturing (decorticate or decerebrate rigidity)
\item Loss of corneal, oculocephalic, or gag reflexes
\item Asymmetric motor responses or abnormal eye movements
\item Seizure activity or posturing
\end{itemize}

\section*{How It Develops}
Consciousness requires an intact ascending reticular activating system (ARAS) in the brainstem and functional connections to the cerebral cortex. Coma results from either diffuse bilateral cortical injury (hypoxia, toxins, metabolic derangement) or direct damage to the ARAS (stroke, hemorrhage, compression). Metabolic causes interfere with neuronal energy metabolism. Structural causes disrupt neural pathways directly.

\section*{Causes}
\begin{itemize}
\item Traumatic brain injury (most common cause in young adults)
\item Stroke (ischemic or hemorrhagic)
\item Hypoxic-ischemic encephalopathy (cardiac arrest, near-drowning)
\item Drug overdose (opioids, benzodiazepines, barbiturates)
\item Metabolic disorders (hypoglycemia, hepatic encephalopathy, uremia)
\item CNS infection (meningitis, encephalitis)
\item Brain tumor with herniation
\item Seizure (status epilepticus, post-ictal state)
\item Carbon monoxide poisoning
\item Hypothermia or hyperthermia
\item Electrolyte disturbances
\end{itemize}

\section*{Related Symptoms}
\begin{itemize}
\item Loss of consciousness
\item Seizures
\item Syncope (transient loss of consciousness)
\item Delirium (acute confusion)
\item Brain herniation
\item Respiratory failure
\item Hypotension or shock
\end{itemize}

\section*{Medical Evaluation}
\begin{itemize}
\item Glasgow Coma Scale (GCS) assessment
\item Rapid blood glucose measurement
\item Complete blood count, electrolytes, renal and liver function
\item Arterial blood gas analysis
\item Toxicology screen for drug overdose
\item Non-contrast CT head for structural causes
\item MRI brain for more detailed assessment
\item Electroencephalogram (EEG) to detect seizures
\end{itemize}

\section*{Treatment Options}
\begin{itemize}
 \item Immediate airway, breathing, and circulation support; advanced airway management is based on protective reflexes, breathing, oxygenation, and the overall clinical situation
\item Correct metabolic derangements (dextrose, naloxone for opioids)
\item Control intracranial pressure if elevated
\item Treat underlying cause (antibiotics, surgery)
\item Antiepileptic drugs for status epilepticus
 \item Temperature control and intensive care after selected cardiac-arrest or brain-injury cases
\item Intensive care monitoring and supportive care
\end{itemize}
\section*{Self-Care and Home Management}
\begin{itemize}
\item This condition requires emergency medical care, not self-management.
\item If a person suddenly becomes unresponsive, call emergency services immediately.
\item Do not give food, drink, or medications to an unconscious person.
 \item If the person is breathing normally, place them in the recovery position unless a spinal injury is suspected; avoid moving them unnecessarily.
 \item Check for normal breathing and follow the emergency dispatcher's instructions; begin CPR if they are not breathing normally.
\end{itemize}

\section*{References}
\begin{thebibliography}{99}
\bibitem{coma:ref1} Wijdicks EFM. ``The Comatose Patient.'' 2nd ed. Oxford University Press, 2014.
\bibitem{coma:ref2} Plum F, Posner JB. ``The Diagnosis of Stupor and Coma.'' 4th ed. Oxford University Press, 2007.
\bibitem{coma:ref3} Giacino JT, Katz DI, Schiff ND, et al. ``Practice Guideline: Disorders of Consciousness.'' Neurology. 2018;91(10):450--460.
\bibitem{coma:ref4} Young GB. ``Coma.'' Ann Neurol. 1999;45(4):417--424.
\bibitem{coma:ref5} Stevens RD, Bhardwaj A. ``Approach to the Comatose Patient.'' Crit Care Med. 2006;34(1):31--41.
\bibitem{coma:ref6} Edlow JA, Rabinstein AA, Traub SJ, et al. ``Diagnosis of Altered Mental Status in the Emergency Department.'' Ann Emerg Med. 2013;61(5):569--579.
\end{thebibliography}
